\section{A degree-706 equation in ten positive integers}
\label{dy:section}

The first construction preserves the ten positive unknowns of the
scaled-norm candidate in the pinned repository, while reducing its exact
total degree from $816$ to $706$.  Its division-free circuit uses $133$
binary arithmetic operations, compared with $124$ for that parent.
This is a degree--operation tradeoff, rather than a claim that one measure
of simplicity captures every consideration.  The coefficient tuple and
its calibration remain those described in Section~\ref{sec:calibration}.
The immediate mathematical parent is~\cite{parent816}; the source interface
and the complete fixed coefficient binding are~\cite{source61,source44}.

The main change is to make all four rational radicands units at the prime
$2$.  Distinct powers of $2$ then separate their rational square classes.
This removes a size-order condition on the strong Pell root and permits
one factor to be removed from its ordinate.  Three centers in the carrier
and a quadratic transport majorant account for the remaining reduction.
We prove the new algebra and positivity arguments explicitly.  The
inherited compiler/calibration and the stated Pell index lemmas retain
precisely the status recorded in Section~\ref{sec:calibration}.

\subsection{The polynomial and its fixed coefficients}
\label{dy:definition}

Write the four fixed integer coefficients as
\[
 \beta=2^{590069}-1,\qquad K_0,\qquad M_C,\qquad M_F.
\]
The complete binary payloads are supplied with the circuit and authenticated
by their compressed and uncompressed SHA-256 digests.  Their verified
bounds are
\begin{equation}
  K_0>0,\qquad 0<M_C<\beta<M_F<2\beta.
  \label{dy:ports}
\end{equation}
Here $M_F$ is already the shifted high mask; it must not be shifted a
second time.  None of these four constants is a quantified unknown.

The ten unknowns are
\begin{equation}
 (j,F,\alpha,i,T,w,Y,n,\eta,\zeta)\in\Pos^{10},
 \qquad \Pos=\{1,2,3,\ldots\}.
 \label{dy:variables}
\end{equation}
Every register in the following definition is an integer polynomial in
these unknowns and the four fixed constants.  Define the outer registers
by
\begin{align}
 J&=2j+1,& q&=\beta J+1,& C&=q-F-\alpha, \notag\\
 X&=wq,& a&=Y(X+1),& k&=\eta+\zeta, \notag\\
 c&=kY+\eta,& \mu&=q^2+1,& \lambda&=2\mu.
 \label{dy:outer}
\end{align}
The scaled main and first Pell parameters are
\begin{align}
 A&=\lambda(a+2)+1,& H&=4A-5,& \Delta&=A^2-1, \notag\\
 a_-&=A-2,& V_0&=X\lambda^2Y^2,&
 D_\tau&=V_0(V_0+1)k^2+1.
 \label{dy:parameters}
\end{align}
The packed index, its excess and the transport numerator are
\begin{align}
 G&=(q-1)(q-F)+\alpha, \notag\\
 R&=G(q^2-1)+(M_C+qM_F)J, \notag\\
 e&=k-R-1, & B_t&=(K_0+w)C-F.
 \label{dy:index}
\end{align}
For the strong and auxiliary Pell registers put
\begin{align}
 I&=ic^2,& D_f&=1+\Delta I^2,& Q&=\Delta^2I^2,
       &\delta&=Q-1, \notag\\
 T'&=T+\alpha,& b&=cT'-R,& V&=D_fb-c, \notag\\
 D_y&=QV^2-1,& D_m&=\Delta c^2+1.
 \label{dy:auxiliary}
\end{align}
In particular, $I$ has no factor $\mu$ other than any accidental
divisibility of its value.  Define the centers and transport weight
\begin{align}
 C_f&=AI,& C_\tau&=V_0k, \notag\\
 L&=(k+F+1)(T'+6),& W&=1+(q-1)L,
 \label{dy:weight}
\end{align}
and the denominator-clearing registers
\begin{align}
 P_1&=q^3(q-1),& M&=P_1\mu^2H, \notag\\
 Z&=P_1\mu^2\bigl[H(n+V+2C_f+4C_\tau)+8(X+a_-c)\bigr]
       \notag\\
  &\hspace{1.2em}+H\bigl[\mu^2((q-1)Y-q^3WB_t)-P_1e\bigr].
 \label{dy:carrier}
\end{align}
For the final norm use the five polynomial arguments
\begin{align}
 s&=\delta Z^2,&
 u_1&=16\delta M^2D_\tau,&
 u_2&=4\delta M^2D_f, \notag\\
 u_3&=M^2D_y,&
 u_4&=64\delta(P_1\mu^2)^2D_m.
 \label{dy:normargs}
\end{align}

Here is a compact fixed norm polynomial in five formal arguments.  Its
letters $a,b,c,d$ are local formal arguments, independent of the similarly
named registers above:
\begin{align}
 \sigma_1&=a+b+c,& \sigma_2&=ab+(a+b)c,& \sigma_3&=abc, \notag\\
 U&=s+d,& P&=U-\sigma_1,& E&=P^2+4(sd-\sigma_2), \notag\\
 G_*&=E^2+16(sdP^2-4\sigma_3U),&
 O_*&=EP-16\sigma_3, \notag\\
 \mathcal F_{16}(s;a,b,c,d)&=G_*^2-64sdO_*^2.
 \label{dy:f16}
\end{align}
The equation delivered here is
\begin{equation}
 \boxed{\;N=\mathcal F_{16}(s;u_1,u_2,u_3,u_4)=0,\quad
         (j,F,\alpha,i,T,w,Y,n,\eta,\zeta)\in\Pos^{10}.\;}
 \label{dy:equation}
\end{equation}
The named registers are abbreviations for substitution, not additional
unknowns.  Equations~\eqref{dy:outer}--\eqref{dy:f16} are consequently a
finite, exact specification of one integer-coefficient polynomial.

\begin{theorem}[Ten-positive equation]
\label{dy:main}
With the fixed coefficient binding and source calibration of
Section~\ref{sec:calibration}, the polynomial $N$ has a zero in
$\Pos^{10}$ if and only if ZFC has a finite syntactic refutation.
The reduction is ordinary mathematics formalizable in ZFC.  Thus the
source calibration gives
\[
 \mathrm{ZFC}\vdash
 \bigl[(\forall\mathbf x\in\Pos^{10})\,N(\mathbf x)\ne0\bigr]
       \longleftrightarrow\Con(\mathrm{ZFC}).
\]
The exact total degree of $N$ is $706$.  The supplied literal circuit has
$133$ binary operations: $79$ multiplications, $33$ additions and $21$
subtractions, including operations with small numerals.
\end{theorem}

\subsection{The norm identity}
\label{dy:normproof}

We include the conjugate-product calculation because it specifies exactly
what the single polynomial equation asserts.

\begin{lemma}[Sixteen conjugates]
\label{dy:normlemma}
In the polynomial ring over $\Z$,
\begin{equation}
 \mathcal F_{16}(x^2;r_1^2,r_2^2,r_3^2,r_4^2)
 =\prod_{\varepsilon\in\{1,-1\}^4}
       (x-\varepsilon_1r_1-\varepsilon_2r_2
          -\varepsilon_3r_3-\varepsilon_4r_4).
 \label{dy:normidentity}
\end{equation}
Moreover, $\mathcal F_{16}$ is homogeneous of degree $8$ in its five
arguments and has leading coefficient $1$ in $s$.
\end{lemma}

\begin{proof}
Put $a=r_1^2$, $b=r_2^2$, $c=r_3^2$ and $d=r_4^2$.
Pairing the signs of $r_1$ and $r_2$, then the sign of $r_3$, gives the
three-root identity
\[
 \prod_{\varepsilon_1,\varepsilon_2,\varepsilon_3}
   (y-\varepsilon_1r_1-\varepsilon_2r_2-\varepsilon_3r_3)
 =\bigl[(y^2-\sigma_1)^2-4\sigma_2\bigr]^2
       -64\sigma_3y^2.
\]
For example the first pairing is
$(y^2-a-b)^2-4ab$; substituting $y-r_3$ and $y+r_3$ and multiplying yields
the displayed expression.

Now fix the sign $\varepsilon$ of $r_4$ and set
$y=x-\varepsilon r_4$, $s=x^2$ and $r=2xr_4$.  Then
$y^2=U-\varepsilon r$ and $r^2=4sd$.  Consequently
\[
 (y^2-\sigma_1)^2-4\sigma_2=E-2\varepsilon rP.
\]
The three-root product becomes
\[
 (E-2\varepsilon rP)^2-64\sigma_3(U-\varepsilon r)
       =G_*-4\varepsilon rO_*.
\]
Multiplication over the two values of $\varepsilon$ gives
$G_*^2-16r^2O_*^2=G_*^2-64sdO_*^2$, as claimed.  The product identity also
shows the homogeneity and the coefficient of $s^8$.
\end{proof}

For later use define positive real roots, whenever the positive input is
fixed, by
\begin{equation}
 \tau=\sqrt{D_\tau},\qquad f=\sqrt{D_f},\qquad
 y=\sqrt{D_y/\delta},\qquad D=\sqrt{D_m},\qquad d=D/H.
 \label{dy:roots}
\end{equation}
Their existence for \emph{every} positive input is checked below, before
any equation or transport sign is assumed.  Lemma~\ref{dy:normlemma}
then yields
\begin{equation}
 N=\delta^8M^{16}
    \prod_{\varepsilon\in\{1,-1\}^4}
    \left(\frac ZM-4\varepsilon_1\tau-2\varepsilon_2f
                       -\varepsilon_3y-8\varepsilon_4d\right).
 \label{dy:factorization}
\end{equation}
Direct division of the definition of $Z$ gives the explanatory identity
\begin{align}
 \frac ZM={}&n+V+2C_f+4C_\tau+\frac{8(X+a_-c)}H
       +\frac{Y}{q^3}-\frac{WB_t}{q-1}-\frac e{\mu^2}.
 \label{dy:carrieridentity}
\end{align}
Only the polynomial definition~\eqref{dy:carrier}, not division, is used
in evaluating $N$.

\subsection{Square classes separated by powers of two}
\label{dy:separation}

\begin{lemma}[Dyadic separation]
\label{dy:dyadic}
Let $r_1,\ldots,r_m$ be positive real numbers with rational squares and
$v_2(r_i^2)=0$.  Let $h_1,\ldots,h_m$ be distinct integers.  If, for some
signs $\epsilon_i$, the sum
$\sum_i\epsilon_i2^{h_i}r_i$ is rational, then all the $r_i$ are rational.
\end{lemma}

\begin{proof}
The extension $K=\Q(r_1,\ldots,r_m)$ is a finite multiquadratic Galois
extension.  Its group sends each $r_i$ to $\pm r_i$, defining a character
$\chi_i$.  Two roots have the same character exactly when their ratio is
rational; a character is trivial exactly when its root is rational.
If some nontrivial character $\chi$ occurs, apply its projection
\[
  \frac1{|\operatorname{Gal}(K/\Q)|}
  \sum_{g\in\operatorname{Gal}(K/\Q)}\chi(g)g
\]
to the rational relation.  It yields
\[
 \sum_{\chi_i=\chi}\epsilon_i2^{h_i}r_i=0.
\]
Divide by any root in this nonempty group.  The ratios are rational and
have $2$-adic valuation zero, because their squares have valuation zero.
The resulting rational summands have pairwise distinct valuations $h_i$.
The term of least valuation cannot cancel.  Equivalently, clear the odd
denominators and divide by the least power of two; one summand is odd and
all the others are even.  This contradiction excludes every nontrivial
character.
\end{proof}

The lemma uses valuations of rational ratios only.  It does not require
an extension of the rational $2$-adic valuation to arbitrary irrational
real numbers, and it imposes no relative-size assumption on the roots.

\begin{lemma}[All-input positivity and units]
\label{dy:allinput}
For every input in $\Pos^{10}$, $\delta,D_y,H$ are positive and the four
radicands in~\eqref{dy:roots} are positive rational numbers of $2$-adic
valuation zero.
\end{lemma}

\begin{proof}
The number $\beta$ is odd, as is $J$, so $q$ is even and $q\ge16$.
Thus $\mu$ is odd, $\lambda$ is even, $A$ is odd, $\Delta$ is even,
$V_0$ is even and $H$ is odd.  Also $k\ge2$, $c\ge3$, $I\ge c^2$ and
$\Delta>\mu^2>1$.  In particular $Q>1$ and $\delta>0$.

If the integer $b=cT'-R$ is zero, then $V=-c$ and $|V|\ge3$.
If $b\ne0$, then
\[
 |V|=|D_fb-c|\ge D_f-c>1,
\]
since $D_f\ge1+\Delta c^4$.  Hence $D_y=QV^2-1>0$ in both cases.
Every numerator and denominator in
\[
 D_\tau,\qquad D_f,\qquad \frac{D_y}{\delta},\qquad
 \frac{D_m}{H^2}
\]
is odd: $D_\tau=1+V_0(V_0+1)k^2$, $D_f=1+\Delta I^2$,
$Q=\Delta^2I^2$ is even, and $D_m=1+\Delta c^2$.
All four rational numbers are therefore units at $2$.
\end{proof}

\begin{lemma}[Adjacent-coefficient descent]
\label{dy:descent}
Let $Q>1$ and $V$ be integers.  If $y\in\Q$ satisfies
\[
 QV^2-(Q-1)y^2=1,
\]
then $y$ is an integer.
\end{lemma}

\begin{proof}
Replace $V,y$ by their absolute values.  Neither can be zero.
Write $y=u/m$ in lowest terms, with $m>0$.  The equation implies
$m^2\mid Q-1$.  When $V=1$, it gives $y=1$.
For $V>1$ define
\[
 V'=(2Q-1)V-2(Q-1)y,\qquad
 y'=(2Q-1)y-2QV.
\]
The first quantity is an integer because $m\mid Q-1$.  The identity
$(2Q-1)^2-4Q(Q-1)=1$ gives
\[
 Q(V')^2-(Q-1)(y')^2=1.
\]
Furthermore
\[
 (Q-1)(y^2-V^2)=V^2-1>0
\]
shows $y>V$, whereas
\[
 (QV)^2-((Q-1)y)^2=QV^2+Q-1>0
\]
shows $QV>(Q-1)y$.  These inequalities imply $-V<V'<V$.
The new equation excludes $V'=0$.  Since $2Q-1\equiv1\pmod m$,
the reduced denominator of $y'$ is still $m$.
Induction on the positive integer $|V|$, using $|V'|<V$, reaches $V=1$.
Thus $m=1$.
\end{proof}

If $N=0$, equation~\eqref{dy:factorization} gives a rational signed
sum of the roots with weights $(1,2,4,8)$ in the order $(y,f,\tau,d)$.
Lemmas~\ref{dy:allinput} and~\ref{dy:dyadic} make all four roots rational.
The roots $\tau,f,D$ have integer squares and hence are integers.
Lemma~\ref{dy:descent} applied to
$QV^2-(Q-1)y^2=1$ makes $y$ an integer as well.  This step is valid even
when $V$ is still negative: the descent can use $|V|$.

\subsection{Soundness: extracting transport and the canonical data}
\label{dy:soundness}

Assume $N=0$ at a positive input, and choose signs from a vanishing factor
of~\eqref{dy:factorization}.  The first congruence is obtained before
transport positivity.  Since $q$ is even,
\[
 \gcd(\mu^2,q^3(q-1))=1,
 \qquad H=4\lambda(a+2)-1\equiv-1\pmod\mu.
\]
Thus $\mu^2$ is coprime to $q^3(q-1)H$.  Integer-valuedness of the
signed roots in~\eqref{dy:carrieridentity} gives
\[
 \frac{8(X+a_-c-\varepsilon_4D)}H+\frac Y{q^3}
           -\frac{WB_t}{q-1}-\frac e{\mu^2}\in\Z.
\]
Clearing the other denominators therefore gives $\mu^2\mid e$.
Put
\begin{equation}
 \rho=\frac e{\mu^2}\in\Z,\qquad
 g=\frac{\varepsilon_4D-X-a_-c}{H},\qquad
 s_0=\frac Y{q^3}>0,\qquad t=\frac{B_t}{q-1}.
 \label{dy:fractionnames}
\end{equation}
At this stage $g,s_0,t$ need not be integers.  The carrier identity is
\begin{align}
 n={}&4(\varepsilon_1\tau-C_\tau)
       +2(\varepsilon_2 f-C_f)+(\varepsilon_3y-V)
       +8g-s_0+Wt+\rho.
 \label{dy:nequation}
\end{align}

Use the standard Pell notation
\[
 \chi_B(r)+\psi_B(r)\sqrt{B^2-1}
       =(B+\sqrt{B^2-1})^r.
\]
The integer roots give, for positive indices,
\begin{align}
 \tau&=\chi_P(n_0),&k&=2\psi_P(n_0),&P&=2V_0+1,\notag\\
 D&=\chi_A(p),&c&=\psi_A(p),\notag\\
 f&=\chi_A(m),&I&=\psi_A(m).
 \label{dy:pellindices}
\end{align}
For the first equation the fundamental solution of
$x^2-V_0(V_0+1)k^2=1$ is $(2V_0+1,2)$: $k=1$ would require a square
strictly between $V_0^2$ and $(V_0+1)^2$.
If $n_0=1$, then $k=2$, $\eta=\zeta=1$ and $c=2Y+1$.
But $\psi_A(1)=1$ and $\psi_A(p)\ge2A>2Y+1$ for $p\ge2$.
Consequently $n_0\ge2$, and
\begin{equation}
 c>k\ge4P>8X\lambda^2Y^2>8q^5.
 \label{dy:earlyc}
\end{equation}

The mask bound~\eqref{dy:ports} gives
\[
 0<(M_C+qM_F)J<2q^2.
\]
If $F\ge q$, then $G\le\alpha$; if $F<q$, then
$G<(q-1)^2+\alpha$.  In either case
\begin{equation}
 R<q^4+(\alpha+2)q^2.
 \label{dy:Ronesided}
\end{equation}
Together with~\eqref{dy:earlyc}, this proves $c(T+\alpha)>R$.
Hence $b\ge1$, $V=D_fb-c>1$ and $y>V$.
No conclusion about $B_t$ has yet been used.

\begin{lemma}[Quadratic transport bound]
\label{dy:transport}
Every positive zero of $N$ satisfies $B_t\ge0$.
\end{lemma}

\begin{proof}
Before any sign restriction on $R$, the definition of $R$ gives
\begin{equation}
 |R|<q^4+Fq^3+(\alpha+2)q^2.
 \label{dy:Rabsolute}
\end{equation}
As $q\ge16$ and $\mu^2>q^4$, it follows that
\begin{equation}
 \frac{|R|}{\mu^2}<F+\alpha+1,\qquad
 |\rho|\le\frac{k+1+|R|}{\mu^2}<k+F+\alpha+1.
 \label{dy:rhoestimate}
\end{equation}
The root centers satisfy the one-sided estimates
\begin{equation}
 \tau-C_\tau<k,\qquad f-C_f<0,\qquad 8g<k.
 \label{dy:rootupper}
\end{equation}
For the first inequality, square the comparison with $(V_0+1)k$.
For the second, $D_f=A^2I^2-I^2+1<A^2I^2$, since $I>1$.
For the third, $D<Ac$ implies $g<2c/H$; now $c<2kY$ and $H>32Y$.
These upper bounds apply to the corresponding signed roots in
\eqref{dy:nequation} as well.

Put $u=y-V$.  Since $y>V>1$,
\begin{align*}
 u&=\frac{V^2-1}{(Q-1)(y+V)}
       <\frac{V}{2(Q-1)}
       <\frac{D_fb}{2(Q-1)}<\frac b\Delta.
\end{align*}
The last inequality is equivalent to
$\Delta D_f=Q+\Delta<2Q-2$, which holds because
$Q=\Delta^2I^2>\Delta+2$.  Also $\Delta>\mu^2$ and
$c<2kY<k\Delta$.  Thus
\begin{equation}
 u<kT'+F+\alpha+1.
 \label{dy:gapupper}
\end{equation}

If the integer $B_t$ were negative, it would be at most $-1$, and
$Wt\le-W/(q-1)<-L$.  Substituting
\eqref{dy:rhoestimate}--\eqref{dy:gapupper} into
\eqref{dy:nequation}, and using
$\varepsilon_3y-V\le y-V$ and $s_0>0$, yields
\[
 n<kT'+6k+2F+2\alpha+2-L<0.
\]
Indeed
\begin{align*}
 L-(kT'+6k+2F+2\alpha+2)
  &=(F+1)T'+4F+4-2\alpha\\
  &=(F+1)T+(F-1)\alpha+4F+4>0.
\end{align*}
This contradicts the stipulated $n\ge1$.
\end{proof}

Now $(K_0+w)C\ge F>0$, so $C\ge1$.  The source outer estimates give
\begin{equation}
 3q<R<q^4-2.
 \label{dy:Router}
\end{equation}
For completeness, write
$G=(q-1)C+q\alpha$ and use $C+\alpha=q-F\le q-1$.
Its lower bound and the positive mask term imply the left inequality;
$G\le q^2-q$ and the mask bound below $2q^2$ imply the right one.

We record how the inherited index arguments apply; this is also a check
that deleting the factor $\mu$ from $I$ changes no hypothesis.
The first Pell congruence is
\[
 k\equiv2n_0\pmod{4V_0},\qquad \mu^2\mid4V_0.
\]
Together with $\mu^2\mid k-R-1$ and~\eqref{dy:Router}, it gives
$2n_0\ge R+1$.  Equations~\eqref{dy:earlyc} and~\eqref{dy:Router}
give $c>2R$; the main Pell equation gives $c>2p$.
The rank lemma requires $c^2\mid I$, which holds exactly as written in
\eqref{dy:auxiliary}.  It gives $c\mid m$ and $m>2p$.
With $S=\Delta I$ the auxiliary data satisfy
\[
 S^2=\Delta(f^2-1)=Q,\qquad c\mid S,\qquad
 V\equiv-c\pmod f,\qquad V\equiv-R\pmod c.
\]
The source index lemma therefore gives $p=R$.
Since $P>A$, monotonicity of Pell ordinates gives $n_0<R$.
The congruence $2n_0\equiv R+1\pmod{\mu^2}$, together with
$R<q^4<\mu^2$, forces $2n_0=R+1$.  The fixed-minus parity lemma then
implies $R\equiv3\pmod4$.

Set $r=(R-1)/2$.  The parent ratio lemma is valid for every integer
scaling $\lambda\ge1$ and hence for $\lambda=2\mu$.  In its notation,
\begin{equation}
 \xi=\frac{(X+1)^{2r}}{X^r},\qquad
 \xi<\frac{\psi_A(2r+1)}{\psi_P(r+1)}
       <\xi\left(1+\frac3a\right)^{2r}.
 \label{dy:ratio}
\end{equation}
The defining interval $kY<c<k(Y+1)$ gives
\begin{equation}
 Y>\frac{X^r}{3},\qquad a>\frac{X^{r+1}}3,\qquad r\ge24.
 \label{dy:Ybound}
\end{equation}
Before the exponent has been determined, clearing the three remaining
fractions gives
\[
 H\mid8q^3(q-1)(X+a_-c-\varepsilon_4D).
\]
Because $H$ is odd, the factor $8$ cancels.  The main projection identities
are
\[
 D-a_-c\equiv2^R\pmod H,\qquad
 D+a_-c\equiv2^{-R}\pmod H.
\]
The parent Archimedean estimate, using~\eqref{dy:Ybound}, is
\begin{equation}
 q^4,2^{R+1}X<X^{r+1}<3a<H.
 \label{dy:archimedean}
\end{equation}
If $\varepsilon_4=-1$, multiplication by $2^R$ would make
$q^3(q-1)(2^RX+1)$ a positive multiple of $H$ smaller than $H$.
This is impossible.  If $\varepsilon_4=+1$, the same estimate bounds
$|q^3(q-1)(X-2^R)|$ strictly below $H$, so divisibility gives
\begin{equation}
 \varepsilon_4=+1,\qquad X=2^R.
 \label{dy:exponent}
\end{equation}
In particular $g$ is an integer.  The remaining partial fractions now give
$q^3\mid Y$ and $q-1\mid B_t$.

The ratio estimate~\eqref{dy:ratio} and the half-binomial lemma yield
\begin{equation}
 Y=\frac12\sum_{h=0}^r\binom{2r}{r+h}X^h.
 \label{dy:halfbinomial}
\end{equation}
To see the uniqueness mechanism, write
\[
 \xi=M_0+\theta,\qquad
 M_0=\sum_{h=0}^r\binom{2r}{r+h}X^h,\qquad0<\theta<\tfrac14.
\]
Here $M_0$ is even because $X=2^{2r+1}$ and the central binomial
coefficient is even.  The quantitative ratio estimate puts $c/k$ in
$(M_0/2,M_0/2+1)$, while its defining interval is $(Y,Y+1)$.
These two integer parts must agree.  The explicit estimates and source
reconstruction are the ones in~\cite[Sections 6.17--6.18]{parent816};
only their arbitrary positive scaling parameter has changed.

Finally keep $J,F,\alpha,w,Y$ and set
$s=Y/q^3$, $z=B_t/(q-1)+1$.
Reconstruct fresh \emph{unscaled} Pell witnesses with main parameter
$a+2$ and first parameter $2XY^2+1$, using
\eqref{dy:exponent}, $R\equiv3\pmod4$ and~\eqref{dy:halfbinomial}.
This is the source reconstruction of~\cite[Section 6.18]{parent816}:
it requires only the recovered canonical data and supplies all fifteen
positive witnesses of the calibrated source system.
Theorem~\texttt{thm:source} of~\cite{source61} gives a finite syntactic
refutation of ZFC.  This proves soundness relative to the calibrated
interface.

\subsection{Completeness and the positive sum coordinate}
\label{dy:completeness}

Suppose ZFC has a finite syntactic refutation.  The calibrated source
provides a canonical positive solution with
\[
 q=B^N\text{ for odd }N>1,\quad R\equiv3\pmod4,\quad
 X=2^R,\quad C>0,
\]
and $Y$ equal to~\eqref{dy:halfbinomial}.  Its odd $J$ gives the positive
integer $j=(J-1)/2$.
The parent ratio argument, valid for arbitrary positive integer scaling,
constructs main and first Pell values for $\lambda=2\mu$:
\[
 c=\psi_A(R),\quad D=\chi_A(R),\quad
 k=2\psi_P((R+1)/2),\quad \tau=\chi_P((R+1)/2).
\]
Set $\eta=c-kY$, $\zeta=k-\eta$; both are positive.
The quantities
\[
 \rho=(k-R-1)/\mu^2,\qquad
 g=(D-X-a_-c)/H,\qquad t=B_t/(q-1)
\]
are integers with $\rho,g\ge1$ and $t\ge0$, as in the parent
completeness proof.

For the strong Pell pair use the smaller index
\begin{equation}
 m=Rc,\qquad I=\psi_A(m),\qquad f=\chi_A(m).
 \label{dy:strongcomplete}
\end{equation}
The multiplication formula and the standard Pell congruence give
\[
 \psi_A(Rc)=c\psi_D(c),\qquad
 \psi_D(c)\equiv cD^{c-1}\pmod{D^2-1}.
\]
Since $c^2\mid D^2-1$, this implies $c\mid\psi_D(c)$ and therefore
$c^2\mid I$.  Thus $i=I/c^2$ is a positive integer.
No divisibility by $\mu$ is needed.

Let $S=\Delta I$ and $Q=S^2$.  The explicit solution in the auxiliary
completeness lemma has
\begin{align}
 V^*&=\frac{\chi_S(R)}S=Q_v(S^2),&y&=\psi_S(R),&v&=(R-1)/2,\notag\\
 T^*&=\frac{V^*+c+Rf^2}{cf^2},&
 Q(V^*)^2-(Q-1)y^2&=1.
 \label{dy:auxcomplete}
\end{align}
The inherited lemma proves that $T^*$ is an integer.  The associated odd
Pell polynomial obeys
\[
 Q_v(z)\ge(4z-3)^v\quad(z\ge1).
\]
Since $v\ge2$ here,
\begin{equation}
 V^*\ge(4Q-3)^v\ge(3Q)^v\ge9Q^2.
 \label{dy:Vlarge}
\end{equation}
As $c<S$ and $f^2=D_f<Q$, it follows that
\[
 T^*>\frac{V^*}{S^3}\ge9S>\alpha.
\]
The last comparison uses the canonical $C>0$, hence
$\alpha<q<S$.  Define $T=T^*-\alpha\ge1$; then $T'=T^*$ and
$V=V^*$ as required.

There is one new completeness issue: the center $C_f=AI$ lies above
$f$, so its contribution to $n$ is negative.  The auxiliary gap has ample
room to absorb it, with an explicit bound.

\begin{lemma}[Canonical auxiliary gap]
\label{dy:gaplower}
For the witnesses just constructed, $u=y-V$ satisfies
\[
 u>2I+Y+1.
\]
\end{lemma}

\begin{proof}
Since $V\ge2$, $Q\ge3$ and $V<y<2V$,
\[
 u=\frac{V^2-1}{(Q-1)(y+V)}
      >\frac V{4Q}
      \ge\frac{9S^2}{4}>3S+1.
\]
For the first strict inequality one may use
$V^2-1\ge\tfrac34V^2$ and
$(Q-1)(y+V)<3QV$.  The second is~\eqref{dy:Vlarge}.
Finally $Y<c<I<S$, so $3S+1>2I+Y+1$.
\end{proof}

Define the integer
\begin{equation}
 n=Wt+\rho+8g+4(\tau-C_\tau)+2(f-C_f)+u-Y/q^3.
 \label{dy:ncomplete}
\end{equation}
The terms $Wt,\rho,g$ are nonnegative or positive, and $\tau>C_\tau$.
Moreover $f>(A-1)I$, so $f-C_f>-I$, and $Y/q^3\le Y$.
Lemma~\ref{dy:gaplower} therefore gives $n>0$, hence $n\in\Pos$.
Substitution into~\eqref{dy:carrieridentity} yields
\[
 Z/M=4\tau+2f+y+8d.
\]
The all-positive sign factor in~\eqref{dy:factorization} vanishes, proving
$N=0$.  This establishes completeness and the semantic part of
Theorem~\ref{dy:main}.

\subsection{A symbolic proof of the exact degree}
\label{dy:degree}

The literal circuit gives the conservative syntactic degree $832$.
That estimate ignores cancellations caused by the three centers.  We
prove the exact total degree using a formal parameter, so no numerical
specialization is needed for the upper bound.

Replace every unknown $x$ by $tx$, keeping the direction variables
algebraically independent, and work over their rational function field
$K$.  Degree below means degree in $t$ at infinity.  In $K((t^{-1}))$
each relevant square root has a branch whose leading term is its center:
this follows by the binomial expansion of $(1+h)^{1/2}$ when $h$ has
negative degree.  For $y$, choose the branch with leading term $V$.
These formal branches need not be positive real branches; the product
over all signs is independent of that choice.

The register and error degrees are
\begin{center}
\begin{tabular}{@{}lr@{}}
\toprule
Quantity & Degree\\
\midrule
$\mu,\lambda$ & $2$\\
$A,H,\Delta$ & $5,5,10$\\
$I,Q,\delta$ & $5,30,30$\\
$D_f,V,M$ & $20,24,13$\\
$C_f,C_\tau$ & $10,9$\\
$\tau-C_\tau$ & $1$\\
$f-C_f$ & $0$\\
$y-V$ & $-6$\\
$d$ & $2$\\
$L,W,B_t$ & $2,3,2$\\
$Z/M-V-2C_f-4C_\tau$ & $4$\\
\bottomrule
\end{tabular}
\end{center}
For example,
\[
 f^2-C_f^2=1-I^2
\]
has degree $10$, while $f+C_f$ has degree $10$.
Similarly
\[
 (y-V)(y+V)=\frac{V^2-1}{Q-1}
\]
gives degree $48-30-24=-6$ for $y-V$.
The degree-$4$ leading term of the remaining carrier is $-L_2(B_t)_2$,
coming from $-WB_t/(q-1)$; its other terms have degree at most $2$.

The sixteen factors in~\eqref{dy:factorization} have the following
nonzero leading terms:
\begin{center}
\begin{tabular}{@{}p{0.47\linewidth}rrl@{}}
\toprule
Signs & Count & Degree & Leading term\\
\midrule
$\varepsilon_3=-1$ & $8$ & $24$ & $2V_{24}$\\
$\varepsilon_3=1$, $\varepsilon_2=-1$ & $4$ & $10$ & $4(C_f)_{10}$\\
$\varepsilon_3=\varepsilon_2=1$, $\varepsilon_1=-1$
    & $2$ & $9$ & $8(C_\tau)_9$\\
$\varepsilon_3=\varepsilon_2=\varepsilon_1=1$
    & $2$ & $4$ & $-L_2(B_t)_2$\\
\bottomrule
\end{tabular}
\end{center}
Subscripts here denote leading homogeneous forms.  Their product has
exact degree
\[
 8\cdot24+4\cdot10+2\cdot9+2\cdot4=258.
\]
The prefactor $\delta^8M^{16}$ has degree
$8\cdot30+16\cdot13=448$.  Thus
\begin{equation}
 \deg N=258+448=706.
 \label{dy:exactdegree}
\end{equation}

For an explicit check that the leading form is nonzero, put
\begin{align*}
 q_0&=2\beta j,& k_0&=\eta+\zeta,& c_0&=k_0Y,\\
 \lambda_0&=2q_0^2,& A_0&=\lambda_0wq_0Y,&\Delta_0&=A_0^2,\\
 I_0&=ic_0^2,& Q_0&=\Delta_0^2I_0^2,\\
 V_*&=-(\Delta_0I_0^2)q_0^3(q_0-F),\\
 C_{f,0}&=A_0I_0,&
 C_{\tau,0}&=wq_0\lambda_0^2Y^2k_0,\\
 M_0&=4q_0^8A_0,& L_0&=(k_0+F)(T+\alpha),&
 B_{t,0}&=w(q_0-F-\alpha).
\end{align*}
The degree-$706$ component is
\begin{equation}
 2^{22}Q_0^8M_0^{16}V_*^8C_{f,0}^4C_{\tau,0}^2
                    (L_0B_{t,0})^2.
 \label{dy:leadingform}
\end{equation}
Every factor is a nonzero polynomial.  More concretely, on the diagonal
where all ten unknowns equal $t$, put $q_0=2\beta$ as a fixed integer.
The coefficient of $t^{706}$ is
\begin{equation}
 9\cdot2^{202}q_0^{366}(q_0-1)^8(q_0-2)^2>0.
 \label{dy:diagonal}
\end{equation}
This explicit coefficient supplies a lower-degree witness for the actual
fixed coefficient tuple.  The three other coefficient ports occur only
in lower-degree terms.

\subsection{Circuit form, verification and limitations}
\label{dy:verification}

The proof uses $W$ as a convenient abbreviation.  The literal circuit
does not materialize it: it distributes
\begin{equation}
 \frac{WB_t}{q-1}=LB_t+\frac{B_t}{q-1}.
 \label{dy:weightsave}
\end{equation}
Thus the circuit subtracts $LB_t$ directly from
$n+V+2C_f+4C_\tau$, and uses $B_t$ alone in the residual fractional
numerator.  This exact identity saves one binary operation and leaves the
polynomial unchanged.  The final count is $79+33+21=133$; all rows are
live and every unknown and coefficient is used.

The supplied build program checks the four compressed and raw coefficient
hashes, their bit lengths and bounds, the circuit topology, liveness and
operation counts, and deterministic regeneration of the JSON circuit.
The audit evaluates the displayed formulas independently at modular
points and computes complete degree-bounded univariate restrictions
modulo two primes.  These restrictions have degree $706$ and the leading
coefficient~\eqref{dy:leadingform}; the diagonal restrictions also match
\eqref{dy:diagonal}.  Additional exact finite cases check the parity and
majorant premises and the lower bound on the canonical auxiliary gap.
These are implementation checks supporting the proof, rather than a
substitute for it.

The reduction is a measurable improvement in degree over the pinned
parent, with a modest increase in circuit length.  It does not claim the
least possible degree, the fewest possible unknowns, or the shortest
possible equation.  It retains the inherited K215/native compiler route
and its stated formalization boundary.  For a formulation over
nonnegative integers, replacing each unknown $x$ by $x+1$ preserves the
exact degree and adds at most ten additions to the literal evaluator.
